% Generated by src/eval/e6ComputationSaved.js analyze -- do not edit, rerun instead.
% Needs in the preamble: \usepackage{pgfplots}
% Set \ESixData to the directory holding e6_scaling.csv before \input-ing this file.
\providecommand{\ESixData}{e6}
\pgfplotstableset{col sep=comma}

\begin{figure}[htbp]
\centering
\begin{tikzpicture}
\begin{axis}[
  width=0.86\linewidth, height=0.6\linewidth,
  xlabel={Robots in the cell $n$}, ylabel={Interaction spaces computed},
  xmin=2, xmax=20, ymin=0,
  grid=major, grid style={gray!25},
  legend cell align=left,
  legend style={font=\footnotesize, at={(0.03,0.97)}, anchor=north west, draw=none, fill=white, fill opacity=0.8, text opacity=1},
]
  \addplot[black, thick, mark=none] table[x=robots, y=allPairs] {\ESixData/e6_scaling.csv};
  \addlegendentry{without the test, $n(n-1)/2$}
  \addplot[blue, thick, mark=*, mark size=1.5pt] table[x=robots, y=admittedPairs] {\ESixData/e6_scaling.csv};
  \addlegendentry{admitted by the test}
\end{axis}
\end{tikzpicture}
\caption{E6: Interaction-space computations with and without the neighbour test, for cells of $n$ robots on a line at the pitch of the five-robot cell (1.24\,m). The pairs to compute grow as $n(n-1)/2$; the pairs the test admits grow roughly linearly, because a robot has a bounded number of neighbours within twice its reach.}
\label{fig:exp_e6}
\end{figure}
